\section{问题二的模型建立与求解}\label{sec:model2}

\subsection{模型建立}

附件2含随机噪声及固定系统偏差。以方式1为绝对参考系，方式2的观测模型为：
\begin{equation}
    \boldsymbol{z}_2(t) = \boldsymbol{p}(t - \Delta\tau) + \begin{pmatrix} \Delta x \\ \Delta y \end{pmatrix} + \boldsymbol{\varepsilon}_2, \quad \boldsymbol{\varepsilon}_2 \sim \mathcal{N}(\boldsymbol{0}, \sigma_2^2 \boldsymbol{I})
\end{equation}

联合估计 $(\Delta\tau, \Delta x, \Delta y)$ 的目标函数：
\begin{equation}
    (\Delta\tau^*, \Delta x^*, \Delta y^*) = \arg\min \sum_{t_k \in \mathcal{T}} \left\| S_1(t_k) - \left[ S_2(t_k - \Delta\tau) + (\Delta x, \Delta y)^\top \right] \right\|^2
    \label{eq:p2_joint}
\end{equation}

求解采用 Nelder-Mead 法（初始值由粗搜获得）。

\subsection{卡尔曼融合}

对齐后采用\textbf{常加速度模型}异步卡尔曼滤波。状态向量：
\begin{equation}
    \boldsymbol{x}_k = (p_x, p_y, v_x, v_y, a_x, a_y)^\top
\end{equation}

状态转移矩阵 $\boldsymbol{F}(\Delta t)$ 为标准常加速度形式；过程噪声 $\boldsymbol{Q}$ 由加速度抖动强度 $q_a$ 参数化。两路观测按时间戳排序后异步更新，观测矩阵 $\boldsymbol{H} = [\boldsymbol{I}_2, \boldsymbol{0}_{2\times 4}]$。

方式2的观测在更新前先校正偏差：$\tilde{\boldsymbol{z}}_2 = \boldsymbol{z}_2 - (\Delta x^*, \Delta y^*)^\top$。

\subsection{结果与分析}

\begin{table}[H]
    \centering
    \caption{问题二求解结果}
    \label{tab:result2}
    \begin{tabular}{@{}lrl@{}}
        \toprule
        指标 & 数值 & 说明 \\
        \midrule
        $\Delta\tau^*$ & $-48.4413$\,s & 时间偏差 \\
        $\Delta x^*$ & $-3.4353$\,m & 系统偏差（X 方向）\\
        $\Delta y^*$ & $1.8061$\,m & 系统偏差（Y 方向）\\
        对齐 RMSE & $1.5646$\,m & 含噪声合理值 \\
        估计噪声 $\hat\sigma$ & $1.1062$\,m & 与 EDA 吻合 \\
        10\,Hz 轨迹点数 & 8513 & \\
        \bottomrule
    \end{tabular}
\end{table}

联合估计避免了"先估 $\Delta\tau$ 再估 $(\Delta x, \Delta y)$"的误差累积效应。KF 融合利用了两路互补的时间分辨率（4\,Hz + 5\,Hz $\to$ 等效 9\,Hz 采样密度），输出 10\,Hz 平滑轨迹。本文采用的卡尔曼滤波方法基于经典理论框架\cite{welch2006}，并结合多传感器融合定位的最新研究成果\cite{hall1997}。
